%%%%%%%%%%%%%%%%%%%%%%%%%%%%%%%%%%%%%%%%%%%%%%%%%%%%%%%%%%%%%%%%%%%%%
%  GSoC 2026 Proposal — robstatpy: Python Wrappers for RobStatTM   %
%  Version: v4 (working copy; source: proposal_v2.tex)             %
%  Author : Aakarsh Gupta                                          %
%%%%%%%%%%%%%%%%%%%%%%%%%%%%%%%%%%%%%%%%%%%%%%%%%%%%%%%%%%%%%%%%%%%%%

\documentclass[11pt, a4paper]{article}

%--------------------------------------------------------------------
% Packages
%--------------------------------------------------------------------
\usepackage[margin=1in]{geometry}
\usepackage{graphicx}
\usepackage{xcolor}
\usepackage{hyperref}
\usepackage{titlesec}
\usepackage{enumitem}
\usepackage{booktabs}
\usepackage{array}
\usepackage{tabularx}
\usepackage{fancyhdr}
\usepackage{amsmath}
\usepackage{amssymb}
\usepackage{listings}
\usepackage{microtype}
\usepackage{parskip}
\usepackage{mdframed}
\usepackage[T1]{fontenc}
\usepackage{lmodern}

%--------------------------------------------------------------------
% Colors
%--------------------------------------------------------------------
\definecolor{uwpurple}{HTML}{4B2E83}
\definecolor{uwgold}{HTML}{85754D}
\definecolor{accentblue}{HTML}{1A73E8}
\definecolor{lightgray}{HTML}{F5F5F5}
\definecolor{darkgray}{HTML}{2E2E2E}
\definecolor{gsocred}{HTML}{DB4437}
\definecolor{softyellow}{HTML}{FFF8E1}
\definecolor{yellowborder}{HTML}{F9A825}

%--------------------------------------------------------------------
% Hyperref
%--------------------------------------------------------------------
\hypersetup{
  colorlinks  = true,
  linkcolor   = uwpurple,
  urlcolor    = accentblue,
  citecolor   = accentblue,
  pdftitle    = {GSoC 2026: robstatpy, Aakarsh Gupta},
  pdfauthor   = {Aakarsh Gupta},
}

%--------------------------------------------------------------------
% Section styling
%--------------------------------------------------------------------
\titleformat{\section}
  {\large\bfseries\color{uwpurple}}
  {\thesection.}{0.6em}{}
  [\vspace{-0.3em}\textcolor{uwgold}{\titlerule[1.2pt]}]

\titleformat{\subsection}
  {\normalsize\bfseries\color{black}}
  {\thesubsection.}{0.5em}{}

\titlespacing{\section}{0pt}{1.6em}{0.7em}
\titlespacing{\subsection}{0pt}{0.9em}{0.3em}

%--------------------------------------------------------------------
% Header / Footer
%--------------------------------------------------------------------
\pagestyle{fancy}
\fancyhf{}
\renewcommand{\headrulewidth}{0.4pt}
\renewcommand{\headrule}{\hbox to\headwidth{%
  \color{uwpurple}\leaders\hrule height \headrulewidth\hfill}}
\fancyhead[L]{\small\color{uwpurple}\textbf{GSoC 2026}
              \textcolor{black}{$\mid$ \texttt{robstatpy}}}
\fancyhead[R]{\small\color{black}Aakarsh Gupta}
\fancyfoot[C]{\small\color{black}\thepage}

%--------------------------------------------------------------------
% Code listing style
%--------------------------------------------------------------------
\lstdefinestyle{python}{
  language        = Python,
  basicstyle      = \ttfamily\footnotesize,
  keywordstyle    = \bfseries\color{accentblue},
  commentstyle    = \itshape\color{darkgray},
  stringstyle     = \color{gsocred},
  numberstyle     = \tiny\color{darkgray},
  numbers         = left,
  numbersep       = 8pt,
  frame           = single,
  framerule       = 0.4pt,
  rulecolor       = \color{uwpurple!40},
  backgroundcolor = \color{lightgray},
  breaklines      = true,
  showstringspaces= false,
  tabsize         = 4,
}
\lstdefinestyle{rstyle}{
  language        = R,
  basicstyle      = \ttfamily\footnotesize,
  keywordstyle    = \bfseries\color{accentblue},
  commentstyle    = \itshape\color{darkgray},
  stringstyle     = \color{gsocred},
  numberstyle     = \tiny\color{darkgray},
  numbers         = left,
  numbersep       = 8pt,
  frame           = single,
  framerule       = 0.4pt,
  rulecolor       = \color{uwpurple!40},
  backgroundcolor = \color{lightgray},
  breaklines      = true,
  showstringspaces= false,
  tabsize         = 4,
}
\lstset{style=python}

%--------------------------------------------------------------------
% Custom environments
%--------------------------------------------------------------------
\newmdenv[
  backgroundcolor  = softyellow,
  linecolor        = yellowborder,
  linewidth        = 1pt,
  innerleftmargin  = 10pt,
  innerrightmargin = 10pt,
  innertopmargin   = 7pt,
  innerbottommargin= 7pt,
  skipabove        = 4pt,
  skipbelow        = 4pt,
]{todobox}

%====================================================================
%  DOCUMENT
%====================================================================
\begin{document}

%--------------------------------------------------------------------
% TITLE BLOCK
%--------------------------------------------------------------------
\thispagestyle{empty}
\begin{center}
  {\LARGE\bfseries\color{uwpurple}
    \texttt{robstatpy}: Python Wrappers for\\[5pt]
    Modern Robust Statistical Estimation
  }\\[10pt]
  {\large\color{black}
    Google Summer of Code 2026, Project Proposal
  }\\[5pt]
  {\normalsize\color{black}
    Organization: \textbf{R Project for Statistical Computing}
    \quad $\cdot$ \quad
    Submitted: 
  }
\end{center}

\vspace{0.5em}
\noindent\textcolor{uwpurple}{\rule{\linewidth}{2pt}}
\vspace{0.3em}

%--------------------------------------------------------------------
% SECTION 1 — PROJECT TITLE & SHORT DESCRIPTION
%--------------------------------------------------------------------
\section{Project Title and Short Description}

\subsection*{Title}

\textbf{\texttt{robstatpy}: Python Wrappers for the \texttt{RobStatTM}
Robust Statistics R~Package}

\subsection*{Short Description}

The \texttt{RobStatTM} R package supports the Maronna, Martin, Yohai and
Salibian-Barrera (2019) 2nd edition of \emph{Robust Statistics: Theory and
Methods (with R)}. As such, \texttt{RobStatTM} implements state-of-the-art
robust estimators, including bias-robust MM-regression, Rocke's S-estimator
for robust covariance matrices in high-dimensional data, robust PCA via M-scale minimization, robust model
selection via RFPE, and Bianco--Yohai robust logistic regression.
While recent Python packages cover some robust methods, the full range of
\texttt{RobStatTM}'s modern robust methods, fully explained in
\emph{Robust Statistics: Theory and Methods}, remains unavailable in Python.
This project will produce \textbf{\texttt{robstatpy}}, a Python library
exposing the target \texttt{RobStatTM} function set for this project
(Section~4) through lightweight
\texttt{rpy2} wrappers whose numerical outputs reproduce the R originals exactly.
Every wrapper will ship with Python-style documentation modelled on the
\texttt{RobStatTM} R man pages, a \texttt{pytest} validation suite, and
reproducible Jupyter notebook examples.

%--------------------------------------------------------------------
% SECTION 2 — STUDENT INFORMATION
%--------------------------------------------------------------------
\section{Student Information}

\subsection{Contact Details}

\renewcommand{\arraystretch}{1.25}
\begin{tabularx}{\linewidth}{@{} >{\bfseries}l X @{}}
  \toprule
  Field            & Details \\
  \midrule
  Full Name        & Aakarsh Gupta \\
  Mailing Address  & Alta 101, 4735 21st Ave NE, Seattle, WA 98105 \\
  Phone            & +1\,(313)\,663-0680 \\
  University Email & \href{mailto:aakarshg@uw.edu}{\texttt{aakarshg@uw.edu}} \\
  Personal Email   & \href{mailto:aakarshgupta919@gmail.com}
                        {\texttt{aakarshgupta919@gmail.com}} \\
  GitHub           & \href{https://github.com/Aakarsh751}
                        {\texttt{github.com/Aakarsh751}} \\
  Project Repo     & \href{https://github.com/Aakarsh751/robstattm-pyport-AG-prop}
                        {\texttt{github.com/Aakarsh751/robstattm-pyport-AG-prop}} \\
  LinkedIn         & \href{https://linkedin.com/in/aakarsh-gupta-128204218/}
                        {\texttt{linkedin.com/in/aakarsh-gupta-128204218}} \\
  Communication    & Zoom, Google Meet \\
  Time Zone        & Pacific Time (PT,\;UTC$-$7) \\
  \bottomrule
\end{tabularx}
\renewcommand{\arraystretch}{1}

\subsection{Academic Affiliation}

\begin{tabularx}{\linewidth}{@{} l X r @{}}
  \toprule
  \textbf{Institution} & \textbf{Degree} & \textbf{Period} \\
  \midrule
  University of Washington, Seattle
    & B.S.\ in Computational Finance and Risk Management (CFRM)
    & Aug 2025 -- Mar 2027 \\[4pt]
  University of Michigan
    & B.S.\ in Data Science \textit{(Dean's List; International Undergraduate Scholar)}
    & Aug 2023 -- Aug 2025 \\
  \bottomrule
\end{tabularx}

\vspace{0.6em}
\noindent
I am currently a \textbf{junior} with an expected graduation of \textbf{March~2027}.
I began my undergraduate studies in Data Science at the University of Michigan,
where I was placed on the Dean's List and was awarded the International Undergraduate
Scholarship.
In August~2025, I transferred to the University of Washington to pursue a
B.S.\ in Computational Finance and Risk Management, a programme that sits at the
intersection of mathematical statistics, numerical computation, and financial
theory, which is exactly the disciplinary blend required to bridge rigorous
robust-statistics research and the Python data-science ecosystem.

\subsection{Relevant Coursework}

The following UW courses form the academic backbone of this project:

\vspace{0.3em}
\begin{tabularx}{\linewidth}{@{} l X @{}}
  \toprule
  \textbf{Course} & \textbf{Title and Relevance} \\
  \midrule
  \multicolumn{2}{@{}l}{\textit{Completed}} \\[2pt]
  CFRM 425 & \textit{R Programming for Quantitative Finance}:
             hands-on R package development and financial modelling in R \\
  CFRM 410 & \textit{Probability \& Statistics for Computational Finance}:
             estimation theory, asymptotic methods, sampling distributions \\
  CFRM 405 & \textit{Mathematical Methods for Quantitative Finance}:
             optimisation, convexity, and numerical root-finding \\
  AMATH 482 & \textit{Computational Methods for Data Analysis}:
              dimensionality reduction, PCA, and spectral methods in Python \\
  AMATH 352 & \textit{Applied Linear Algebra \& Numerical Analysis}:
              matrix decompositions and numerical stability \\[4pt]
  \multicolumn{2}{@{}l}{\textit{Spring 2026 (in progress)}} \\[2pt]
  CFRM 421 & \textit{Machine Learning with Application to Finance}:
             Python-based ML methods applied to financial data, directly
             mirroring the use-case of \texttt{robstatpy} \\
  CFRM 586 & \textit{Financial Time Series Forecasting Methods}:
             modelling heavy-tailed financial returns, a primary application
             domain for robust estimation \\
  \bottomrule
\end{tabularx}

\subsection{Technical Skills}

\begin{itemize}[leftmargin=1.5em, itemsep=3pt]
  \item \textbf{Programming languages:} Python, R, SQL, C++
  \item \textbf{Python ecosystem:}
        NumPy, pandas, scikit-learn, matplotlib, PyTorch, TensorFlow,
        Selenium; REST API development with FastAPI and Streamlit
  \item \textbf{R ecosystem:}
        \texttt{ggplot2}, \texttt{tidyr}, \texttt{dplyr},
        \texttt{RobStatTM}, \texttt{robustbase};
        DataCamp-certified in R~Programming, Data~Manipulation with dplyr,
        and Data~Visualisation with ggplot2
  \item \textbf{Development tools:}
        Git/GitHub, Jupyter, RStudio, Docker, Anaconda
  \item \textbf{Finance:}
        Bloomberg Finance Fundamentals; quantitative equity and
        crypto-market analysis; portfolio risk modelling
\end{itemize}

\subsection{Relevant Experience and Projects}

\begin{description}[style=nextline, leftmargin=0pt, itemsep=8pt]

  \item[\normalfont\textbf{IT Risk Management and Analyst Intern,
        Virginia Department of Motor Vehicles}
        ]
  \begin{sloppypar}
  Spearheaded a proof-of-concept initiative to evaluate AI and NLP for
  analysing call-centre interactions, assessing customer sentiment,
  call-resolution efficiency, agent performance, and fraud-related risk
  indicators.
  Built FastAPI back-ends and Streamlit front-ends to deploy and interact
  with large language models, enabling real-time inference and operational
  decision support.
  Designed and automated intelligent workflows using n8n to streamline
  repetitive tasks, and processed unstructured speech data to produce
  cost-benefit analyses guiding in-house versus outsourced AI deployment
  decisions.
  This role directly exercised the API-design, packaging, and data-pipeline
  skills that are central to \texttt{robstatpy}.
  \end{sloppypar}

  \item[\normalfont\textbf{Undergraduate Research Volunteer,
        University of Michigan -- Imitation Learning}
        ]
  \begin{sloppypar}
  Conducted research on Imitation Learning under faculty guidance, training
  reinforcement learning and human-guided agents with OpenAI Gym.
  Implemented the DAGGER and HG-DAGGER algorithms from the primary literature
  to improve agent performance, adaptability, and learning efficiency, and
  evaluated all models rigorously using PyTorch.
  This experience demonstrates the ability to read formal algorithm
  specifications, translate them into well-tested, reproducible code,
  and verify results against published benchmarks, a skill directly
  applicable to wrapping the MM- and S-estimation routines in
  \texttt{RobStatTM}.
  \end{sloppypar}

  \item[\normalfont\textbf{Deep Learning Algorithm for Asthenopia Detection}
        ]
  \begin{sloppypar}
  Independently conceived and developed a deep-learning image-classification
  pipeline for the early detection of asthenopia (visual fatigue).
  The project was recognised with a \textbf{Gold Award at the Indian Science
  and Engineering Fair} and achieved a top-100 national selection, providing
  evidence of the capacity to drive an independent research project from
  initial conception through to a peer-recognised outcome.
  \end{sloppypar}

\end{description}

\subsection{Why My Background Aligns with This Project}

My training in Computational Finance and Risk Management (CFRM) is
uniquely suited to this project.
Robust statistics is not an abstract topic for me: it is central to the
quantitative methods I am now beginning to study and apply.
In CFRM~410 (Probability and Statistics), I learned estimation
theory and asymptotic analysis that will allow me to understand the underpinnings of robust MM-estimators and
S-Estimators that are covered in Maronna et al.\ (2019).
In CFRM~425 (R Programming), I developed R packages and gained hands-on
familiarity with the object system that \texttt{RobStatTM} uses
internally.
In CFRM~405 (Mathematical Methods), I studied the optimisation
numerical algorithms that drive iterative robust estimators.
AMATH~482 (Computational Methods for Data Analysis) gave me direct
experience implementing PCA and dimensionality reduction in Python, the
exact techniques that \texttt{pcaRobS} robustifies.

This combination means I understand \emph{both}
the R code base and its statistical foundations, and the Python
ecosystem that \texttt{robstatpy} targets.
I am not just a software engineer wrapping an unfamiliar library; I have
the theoretical knowledge to verify correctness, interpret results, and
extend the project into downstream applications such as robust portfolio
construction and risk analysis (RPCRA).

\subsection{Relationship with Mentors}

I have been in active communication with \textbf{Prof.\ R.\ Douglas Martin}
since \textbf{January 2026}, attending weekly meetings to discuss the project
scope, the Python API design, and the technical challenges of bridging R's S3
object system with Python's language model.
These conversations have directly shaped the implementation strategy presented
in this proposal.
I look forward to further engaging with the full mentorship team,
Prof.\ Matias Salibian-Barrera and Brian G.\ Peterson, as the project progresses.

\subsection{Mentors}

\renewcommand{\arraystretch}{1.2}
\begin{tabularx}{\linewidth}{@{} l l X @{}}
  \toprule
  \textbf{Name} & \textbf{Role} & \textbf{Contact} \\
  \midrule
  Prof.\ R.\ Douglas Martin
    & Primary mentor
    & \href{mailto:martinrd3d@gmail.com}{\texttt{martinrd3d@gmail.com}} \\
  Prof.\ Matias Salibian-Barrera
    & Co-mentor
    & \href{mailto:matias@stat.ubc.ca}{\texttt{matias@stat.ubc.ca}} \\
  Brian G.\ Peterson
    & Co-mentor
    & \href{mailto:brian@braverock.com}{\texttt{brian@braverock.com}} \\
  \bottomrule
\end{tabularx}
\renewcommand{\arraystretch}{1}

\subsection{Time Commitment and Availability}

I am applying for the \textbf{large project size (350~hours)}.
I am fully dedicated and available for the entire duration of the GSoC
coding period (late May through mid-November 2026), with
\textbf{30--35 hours per week} committed to \texttt{robstatpy}
development.
I will continue attending weekly meetings with Prof.\ Martin and will be
responsive to all three mentors via email, Zoom, and GitHub.

%--------------------------------------------------------------------
% SECTION 3 — BENEFITS TO COMMUNITY
%--------------------------------------------------------------------
\section{Benefits to the Community}

The robust-statistics methods described in \texttt{RobStatTM}
including MM-regression, Rocke's S-estimator, robust PCA via M-scale minimization, among others,
are the algorithmic companions to \emph{Robust Statistics: Theory and
Methods} (Maronna, Martin, Yohai \& Salibian-Barrera, 2nd~ed., 2019),
a leading modern textbook in the field.
Today, a Python practitioner who wants to reproduce a textbook example or
apply these methods in their research or applications faces two unappealing choices:
install and learn R, or re-implement algorithms from scratch in Python.

\texttt{robstatpy} removes this barrier by providing a
pip\footnote{\texttt{pip} is the standard Python package installer.
Running \texttt{pip install robstatpy} downloads and installs the
package and all its dependencies automatically.}-installable
Python package that:

\begin{enumerate}[leftmargin=2em]
  \item \textbf{Exposes the target \texttt{RobStatTM} API} (Section~4) through
    \texttt{rpy2} wrappers, accepting NumPy arrays and pandas DataFrames
    and returning native Python objects.
  \item \textbf{Guarantees numerical fidelity:} every output agrees
    with the R original to machine precision, because the R code itself
    executes underneath.
  \item \textbf{Lowers the maintenance burden:} because the wrappers
    delegate to R, any upstream improvement to \texttt{RobStatTM} is
    automatically available in Python without duplicating development
    effort.
  \item \textbf{Reproduces figures and numerical summaries} using
    \texttt{matplotlib} and \texttt{plotnine}\footnote{\texttt{plotnine}
    is a Python implementation of R's \texttt{ggplot2}
    grammar-of-graphics, enabling near-identical plot syntax in Python.
    See \url{https://plotnine.org/} and \url{https://realpython.com/ggplot-python/}.}
    so that Python users can produce visualisations that closely
    mirror the R originals.
  \item \textbf{Provides a comparative study} of wrapper-based versus
    direct-porting strategies, producing benchmarks and guidelines that
    will inform future R-to-Python interoperability projects.
\end{enumerate}

\noindent
\textbf{Who benefits:}
\begin{itemize}[leftmargin=1.5em, itemsep=2pt]
  \item \emph{Data scientists, ML engineers, and Python-only practitioners}
    who need a rich toolkit of modern robust estimation methods without
    learning R. \texttt{robstatpy} provides a comprehensive, self-contained
    Python library: users call familiar Python functions with NumPy arrays
    and pandas DataFrames and never need to write, read, or understand any
    R code.
  \item \emph{Students and researchers} who study the Maronna~et~al.\
    textbook and want to apply the methods in Jupyter rather than RStudio.
\end{itemize}

%--------------------------------------------------------------------
% SECTION 4 — RELATED WORK
% SECTION 4 — RELATED WORK
%--------------------------------------------------------------------
\section{Current \texttt{RobStatTM} Functionality}

The table below lists all \texttt{RobStatTM} R functions recommended
in \textit{Robust Statistics: Theory and Methods} (Maronna, Martin,
Yohai, Salibi\'{a}n-Barrera, 2nd ed.), grouped by chapter. These form
the complete target function set for \texttt{robstatpy}.

\renewcommand{\arraystretch}{1.3}
\begin{tabularx}{\linewidth}{@{} l X r @{}}
  \toprule
  \textbf{R Function} & \textbf{Description} & \textbf{Book \S} \\
  \midrule
  \multicolumn{3}{@{}l}{\textit{Location \& Scale (Chapter 2)}} \\[1pt]
  \quad\texttt{locScaleM}
    & Location and scale M-estimator (joint estimation)
    & 2.3, 2.7 \\
  \quad\texttt{mScale}
    & Scale M-estimator
    & 2.5, 2.6 \\[5pt]
  \multicolumn{3}{@{}l}{\textit{Linear Regression I (Chapter 4)}} \\[1pt]
  \quad\texttt{lmrobM}
    & M-estimator for regression with designed experiments
    & 4.4 \\
  \quad\texttt{rob.linear.test}
    & Robust likelihood-ratio test for linear hypotheses
    & 4.7 \\[5pt]
  \multicolumn{3}{@{}l}{\textit{Linear Regression II (Chapter 5)}} \\[1pt]
  \quad\texttt{lmrobdetMM}
    & MM-estimator with optimal $\rho$; primary regression recommendation
    & 5.3, 5.9 \\
  \quad\texttt{pyinit}
    & Highly robust Pe\~{n}a-Yohai initial estimator for \texttt{lmrobdetMM}
    & 5.7 \\
  \quad\texttt{step.lmrobdet}
    & Robust stepwise model selection via RFPE criterion
    & 5.6 \\
  \quad\texttt{lmrobdetDCML}
    & Distance-constrained MLE; boosts efficiency over MM
    & 5.9 \\
  \quad\texttt{pense}$^{\ast}$
    & Robust elastic net S estimator (MM-lasso as a special case)
    & 5.1 \\[5pt]
  \multicolumn{3}{@{}l}{\textit{Multivariate Analysis (Chapter 6)}} \\[1pt]
  \quad\texttt{covRobMM}
    & MM robust covariance estimator; recommended for $p < 10$
    & 6.5 \\
  \quad\texttt{covRobRocke}
    & Rocke S-type robust covariance estimator; recommended for $p \geq 10$
    & 6.4.2, 6.4.4 \\
  \quad\texttt{KurtSDNew}
    & Pe\~{n}a-Prieto highly robust initial location and covariance matrix estimator
    & 6.9.2 \\
  \quad\texttt{pcaRobS}
    & Robust PCA based on a robust scale estimator
    & 6.11.2 \\
  \quad\texttt{GSE}$^{\ddagger}$
    & Generalized S-estimator for covariance with missing data
    & 6.12.2 \\
  \quad\texttt{TSGS}$^{\ddagger}$
    & Two-step GSE covariance estimator for cell-wise outliers
    & 6.13 \\[5pt]
  \multicolumn{3}{@{}l}{\textit{Generalized Linear Models (Chapter 7)}} \\[1pt]
  \quad\texttt{BYlogreg}$^{\dagger}$
    & M-estimator for logistic regression
    & 7.2 \\
  \quad\texttt{WBYlogreg}$^{\dagger}$
    & Redescending M-estimator for logistic regression (primary GLM recommendation)
    & 7.2 \\
  \quad\texttt{WMLlogreg}$^{\dagger}$
    & Weighted maximum-likelihood robust logistic regression
    & 7.2 \\
  \bottomrule
\end{tabularx}
\renewcommand{\arraystretch}{1}

\vspace{0.3em}
\noindent\footnotesize
$^{\ast}$~\texttt{pense} is an R function provided by the \texttt{pense} package.\\[2pt]
$^{\ddagger}$~\texttt{GSE} and \texttt{TSGS} are R functions provided by the
\texttt{GSE} package.\\[2pt]
$^{\dagger}$~Chapter~7 GLM functions appear here as the book-recommended set;
they are \textbf{not} on the implementation timeline (Section~9).
\normalsize

%--------------------------------------------------------------------
% SECTION 5 -- PROPOSED WORK
%--------------------------------------------------------------------
% SECTION 5 -- TEST TASKS AND VALIDATION
%--------------------------------------------------------------------
\section{GSoC Test Tasks and Validation}

As part of the \texttt{robstatpy} GSoC 2026 application, two required
test tasks were completed: implementing \texttt{rpy2} wrappers for
\texttt{locScaleM} and \texttt{scaleM}, and validating their outputs
against direct R calls. All 14 automated checks pass, confirming
bit-identical numerical agreement between the Python wrappers and
native R.

\noindent\textbf{Test task notebook:}
\href{https://github.com/Aakarsh751/robstattm-pyport-AG-prop/blob/main/robstattm/python/robstatpy\_comparison\_rpy2.ipynb}%
{\texttt{robstatpy\_comparison\_rpy2.ipynb}}

\noindent\textbf{Full wrapper code and result tables:}
\href{https://github.com/Aakarsh751/robstattm-pyport-AG-prop/blob/main/docs/gsoc2026_proposal/robstatpy_tests.pdf}%
{\texttt{robstatpy\_tests.pdf}}

%--------------------------------------------------------------------
% SECTION 6 -- PROPOSED WORK
%--------------------------------------------------------------------
\section{Proposed Work}

\texttt{robstatpy} exposes the target \texttt{RobStatTM} function set
(Section~4) to Python users via lightweight \texttt{rpy2} wrappers, providing
numerically exact access to each in-scope estimator.
This section describes the wrapper approach and presents code examples
for key functions already implemented.

\subsection{Why Wrappers?}

The \texttt{RobStatTM} package is \textbf{actively maintained} by its
original authors, the same researchers who proved the theoretical
properties of these estimators.
A wrapper-based approach has three decisive advantages over full
re-implementation:

\begin{enumerate}[leftmargin=2em]
  \item \textbf{Exact numerical equivalence.}  Because the R code
    executes underneath, Python users get bit-identical results.
    This is essential for textbook reproducibility and for
    cross-validating future pure-Python ports.
  \item \textbf{Zero duplication of maintenance.}  When the R authors
    fix a bug or add an estimator, the Python wrapper inherits the
    improvement immediately.
    A pure-Python re-implementation creates a second code base that
    must be independently maintained and validated.
  \item \textbf{Speed of delivery.}  Wrapping a function requires
    writing conversion glue, tests, and documentation, not
    re-deriving the algorithm.
    This makes it feasible to expose the \emph{full target} API within a
    single GSoC period, rather than a subset.
\end{enumerate}

\noindent
Together, these properties make the wrapper approach the most reliable
path to bringing the in-scope \texttt{RobStatTM} API to Python within a
single GSoC period, without sacrificing any numerical accuracy.

\subsection{\texttt{rpy2} Wrapper Layer}

Every public in-scope \texttt{RobStatTM} function will be wrapped in a Python function
that:
\begin{enumerate}[leftmargin=2em]
  \item Accepts NumPy arrays and/or pandas DataFrames.
  \item Converts inputs to R objects via \texttt{rpy2.robjects} and
    \texttt{numpy2ri}/\texttt{pandas2ri} converters.
  \item Calls the R function inside \texttt{RobStatTM}.
  \item Extracts named elements from the returned R list using
    \texttt{.rx2()}\footnote{In \texttt{rpy2}, \texttt{rx} and \texttt{rx2}
    are \emph{delegators} for R's \texttt{[} and \texttt{[[} operators:
    \texttt{.rx2(\textit{name})} corresponds to \texttt{obj[[\textit{name}]]}
    in R and is the standard way to pull a single named component (e.g.\ a
    coefficient vector or residuals) from an R \texttt{list} or fitted-model
    object. The \texttt{rpy2} vector documentation:
    \url{https://rpy2.github.io/doc/latest/html/vector.html}.}.
  \item Returns a Python dataclass (or named dict) whose fields match
    the R return values.
\end{enumerate}

\noindent
\textbf{Architecture diagram:}

\begin{center}
\small
\setlength{\tabcolsep}{3pt}
\begin{tabular}{ccccc}
  \fbox{\texttt{Python user}} &
  $\xrightarrow{\text{\scriptsize np/pd}}$ &
  \fbox{\texttt{robstatpy}} &
  $\xrightarrow{\text{\scriptsize rpy2}}$ &
  \fbox{\texttt{RobStatTM}} \\[6pt]
  & & $\downarrow$ & & $\downarrow$ \\[4pt]
  & & \fbox{\texttt{dataclass}} &
  $\xleftarrow{\text{\scriptsize \texttt{.rx2()}}}$ &
  \fbox{\texttt{R list}} \\
\end{tabular}
\end{center}
\normalsize

\noindent\textbf{Key design decisions:}

\begin{itemize}[leftmargin=1.5em, itemsep=3pt]
  \item \textbf{Global conversion context.}
    \texttt{rpy2}~3.6+ stores conversion rules in a\\
    \texttt{contextvars.ContextVar}, which is lost across Jupyter's
    async task boundaries.
    We patch this at import time by calling \texttt{set\_conversion}
    with \texttt{default\_converter}, ensuring conversions persist
    across cells.
  \item \textbf{R dot-to-Python underscore mapping.}
    \texttt{importr()} automatically maps\\
    \texttt{lmrobdet.control} to \texttt{lmrobdet\_control}.
    We document both names in every docstring.
  \item \textbf{Deferred R package loading.}
    Calling \texttt{importr("RobStatTM")} at module import time starts R
    immediately; the package design can instead defer \texttt{importr}
    until the first wrapper call so that a bare \texttt{import robstatpy}
    (e.g.\ for type hints or doc builds) avoids initializing R until a
    function is actually used.
\end{itemize}


\subsubsection*{Code Example: \texttt{rpy2} Wrapper for \texttt{lmrobdetMM}}

\begin{lstlisting}[style=python, title={\footnotesize rpy2 wrapper for \texttt{lmrobdetMM} (mopt family, efficiency\,=\,0.95)}]
def lmrobdet_mm(formula, data, family="mopt", efficiency=0.95):
    """Robust MM-regression via RobStatTM::lmrobdetMM().

    Parameters
    ----------
    formula : str
        R-style formula, e.g. 'zinc ~ copper'.
    data : pd.DataFrame
        Data frame (converted to R via pandas2ri).
    family : str
        Psi function family for IRLS refinement.
    efficiency : float
        Target asymptotic efficiency.

    Returns
    -------
    dict with 'coefficients', 'scale', 'residuals',
         'fitted_values', 'rweights', 'r_squared'.
    """
    from rpy2.robjects import pandas2ri
    pandas2ri.activate()
    r_df = ro.conversion.py2rpy(data)
    ctrl = robstattm.lmrobdet_control(
        family=family, efficiency=efficiency)
    fit = robstattm.lmrobdetMM(
        ro.Formula(formula), data=r_df, control=ctrl)
    return {
        "coefficients":  np.array(fit.rx2("coefficients")),
        "scale":         fit.rx2("scale")[0],
        "residuals":     np.array(fit.rx2("residuals")),
        "fitted_values": np.array(fit.rx2("fitted.values")),
        "rweights":      np.array(fit.rx2("rweights")),
        "r_squared":     fit.rx2("r.squared")[0],
    }
\end{lstlisting}



%--------------------------------------------------------------------
% SECTION 7 -- DIRECT PYTHON IMPLEMENTATION (SECONDARY TRACK)
%--------------------------------------------------------------------
\section{What about a Direct Python Implementation?}

A direct Python re-implementation of \texttt{RobStatTM} is \textbf{not the main
priority} of this project. The core deliverable is a complete \texttt{rpy2}
wrapper layer whose outputs match \texttt{RobStatTM} numerically.

Some native Python implementations may be added where appropriate. They provide
an independent code path to the same estimators, so benchmarks can compare
accuracy (e.g.\ agreement with R and with the wrappers), wall-clock cost, and
scaling with problem size side by side with the \texttt{rpy2}-based wrappers,
isolating bridge and conversion effects from the underlying algorithm cost.

% SECTION 8 -- COMPARATIVE BENCHMARKS EVALUATION
%--------------------------------------------------------------------
\section{Comparative Benchmarks Evaluation}
\label{sec:evaluation}

A central deliverable is a rigorous evaluation of the primary
\texttt{rpy2} wrapper path: comparison to direct R calls and, where
relevant, to classical Python estimators, along three dimensions:

\subsection{Metrics}

\begin{enumerate}[leftmargin=2em]
  \item \textbf{Accuracy.}
    For each function, compute $\max_i |y_i^{\text{Py}} - y_i^{\text{R}}|$
    across all return fields for the \texttt{rpy2} wrapper against a direct
    R session call (expected: $0$).
  \item \textbf{Speed.}
    Wall-clock time per call, measured with \texttt{timeit} over 100
    repetitions on synthetic data of sizes $n \in \{100,\,1000,\,10000\}$
    and $p \in \{5,\,20,\,50\}$.
    Breakdown: R computation time vs.\ conversion overhead.
  \item \textbf{Coverage.}
    Number of target functions wrapped and passing all validation checks,
    tracked against the core wrapper set in Section~9
    (e.g., 6/15 at proposal time; 15/15 as the final deliverable for that set).
    Section~4 lists the full 18-function textbook companion set, including
    Chapter~7 (reference only; not on the timeline).
\end{enumerate}

\subsection{Benchmark Protocol}

\begin{itemize}[leftmargin=1.5em, itemsep=3pt]
  \item \textbf{Datasets:}
    Both synthetic and real datasets are used. Synthetic samples are
    drawn from a Gaussian base distribution with $\varepsilon = 5\%$
    contamination --- 5\% of observations are replaced by outliers
    placed $10\sigma$ from the mean --- to stress-test breakdown
    behaviour across sample sizes $n \in \{100,\,1000,\,10000\}$ and
    dimensions $p \in \{5,\,20,\,50\}$.
    Real datasets from the book (mineral, wine, bus) are used to
    reproduce published results.

  \item \textbf{Functions benchmarked:}
    \texttt{lmrobdetMM}, \texttt{covRobMM}, \texttt{covRobRocke},
    \texttt{pcaRobS}.

  \item \textbf{Accuracy comparison:}
    Each \texttt{robstatpy} wrapper output is compared field-by-field
    against a direct R session call (same data, same seed) to confirm
    zero numerical difference. Where a \texttt{scikit-learn} or
    \texttt{statsmodels} equivalent exists, results are also compared
    to quantify the practical difference between robust and
    classical estimators.

  \item \textbf{Speed comparison:}
    Wall-clock time per call measured with \texttt{timeit} over 100
    repetitions. Comparison grid: native R session
    vs.\ \texttt{rpy2} wrapper (showing bridge overhead)
    vs.\ nearest classical baseline (\texttt{scikit-learn} or
    \texttt{statsmodels}, when applicable).

  \item \textbf{Reporting:}
    Results presented in a Jupyter notebook with timing tables,
    accuracy comparison tables, and commentary.
\end{itemize}

%--------------------------------------------------------------------
% SECTION 9 -- IMPLEMENTATION PLAN
%--------------------------------------------------------------------
\section{Implementation Plan and Timeline}

The project follows a \textbf{22-week large-project timeline}
(350~hours), structured around the GSoC 2026 calendar. Testing and
documentation are treated as first-class deliverables and are built
incrementally within each phase rather than deferred.
Section~\ref{sec:resulting-robstatpy} summarises the resulting
\textbf{RobstatPY} features (validation, benchmarks, documentation, tutorials).

\medskip
\noindent\textit{Unless otherwise noted, each weekly milestone includes
field-by-field \texttt{pytest} checks against direct R calls, a NumPy-style
docstring on every new wrapper, and a Sphinx API entry produced with
\texttt{autodoc} (initial stubs are acceptable until the documentation
phase in Weeks~13--16).}

\subsection{Community Bonding Period (May 1 -- May 24)}

\begin{itemize}[leftmargin=1.5em, itemsep=2pt]
  \item Finalise package scaffolding: \texttt{pyproject.toml},
    \texttt{pytest} configuration, CI/CD with GitHub Actions
    (R + Python matrix).
  \item Study \texttt{RobStatTM} vignettes and the 26 example R scripts
    in the project repository to understand expected inputs, outputs,
    and edge cases for each function.
  \item Draft the \texttt{robstatpy} API surface: naming conventions,
    type stubs, and docstring templates modelled on R man pages.
  \item Set up Sphinx + ReadTheDocs skeleton; agree on code review
    cadence and milestone criteria with mentors.
\end{itemize}

\subsection{Phase 1: Univariate Wrappers (Weeks 1--2, May 25 -- Jun 7)}

\begin{tabularx}{\linewidth}{@{} l X @{}}
  \toprule
  \textbf{Week} & \textbf{Deliverables} \\
  \midrule
  1 & \texttt{locScaleM} and \texttt{mScale} \texttt{rpy2} wrappers
    (bisquare, huber, mopt families; eff $\in \{0.85,\,0.90,\,0.95\}$).
    Full \texttt{pytest} suite vs.\ R; document the relationship between
    \texttt{mScale} and \texttt{locScaleM}; numerical cross-checks (same
    inputs and seeds). \\
  2 & \texttt{lmrobdet.control} wrapper exposing all tuning parameters.
    Reproduce Chapter~2 textbook examples in a Jupyter notebook.
    Integration tests for the complete univariate module. \\
  \bottomrule
\end{tabularx}

\subsection{Phase 2: Robust Regression (Weeks 3--6, Jun 8 -- Jul 5)}

\begin{tabularx}{\linewidth}{@{} l X @{}}
  \toprule
  \textbf{Week} & \textbf{Deliverables} \\
  \midrule
  3 & \texttt{lmrobdetMM} wrapper: formula interface, coefficient
    extraction, residuals, robustness weights, scale.
    Pytest suite. \\
  4 & \texttt{lmrobdetMM} diagnostics: robust $R^2$, fitted values,
    summary table mirroring \texttt{summary()} output.
    \texttt{lmrobdetDCML} wrapper + pytest. \\
  5 & \texttt{step.lmrobdet} wrapper for stepwise model selection.
    \texttt{pyinit} and \texttt{rob.linear.test} wrappers.
    Reproduce mineral dataset examples (Figs~5.1--5.7). \\
  6 & \textbf{Midterm milestone:} regression module complete with
    $\geq 90\%$ pytest coverage. Midterm evaluation report. \\
  \bottomrule
\end{tabularx}

\vspace{0.3em}
\noindent\textbf{Midterm Evaluation (Jul 6--10):}
Deliverables: all univariate and regression wrappers scheduled for
Phases~1--2, pytest suite, Chapter~2 and~5 reproduction notebooks,
$\geq 90\%$ test coverage on delivered modules.

\subsection{Phase 3: Robust Covariance and PCA (Weeks 7--10, Jul 13 -- Aug 9)}

\begin{tabularx}{\linewidth}{@{} l X @{}}
  \toprule
  \textbf{Week} & \textbf{Deliverables} \\
  \midrule
  7  & \texttt{covRobMM} wrapper: center, covariance, correlation,
     Mahalanobis distances, weights. Pytest suite. \\
  8  & \texttt{covRobRocke} and \texttt{KurtSDNew} wrappers.
     Reproduce wine dataset examples (Fig~6.3). Pytest suite. \\
  9 & \texttt{pcaRobS} wrapper: eigenvectors, proportion of variance,
     scores. Reproduce bus dataset examples (Fig~6.10). Pytest suite. \\
  10 & Distance--distance plot and scree plot utility functions.
     Documentation pass for all multivariate wrappers.
     End-to-end integration tests for covariance and PCA module. \\
  \bottomrule
\end{tabularx}

\subsection{Phase 4: Benchmarks and Comparative Study (Weeks 11--12, Aug 10 -- Aug 23)}

\begin{tabularx}{\linewidth}{@{} l X @{}}
  \toprule
  \textbf{Week} & \textbf{Deliverables} \\
  \midrule
  11 & Speed benchmarks: native R session vs.\ \texttt{rpy2} wrappers vs.\
     \texttt{scikit-learn}/\texttt{statsmodels} baselines.
     Synthetic contaminated Gaussian data generator. \\
  12 & Accuracy benchmarks: max absolute error across all functions
     and parameter combinations. Tabulate rpy2 bridge overhead.
     Comparative Jupyter notebook with timing tables, accuracy
     results, and commentary. \\
  \bottomrule
\end{tabularx}

\subsection{Phase 5: Documentation, Packaging, and Tutorials (Weeks 13--16, Aug 24 -- Sep 20)}

\begin{tabularx}{\linewidth}{@{} l X @{}}
  \toprule
  \textbf{Week} & \textbf{Deliverables} \\
  \midrule
  13 & Complete Sphinx API documentation for all delivered wrappers.
     Every function has a docstring modelled on the R man page,
     with usage examples. Publish on ReadTheDocs. \\
  14 & Publish \texttt{robstatpy} on PyPI (sdist + wheel).
     Installation guide covering R and \texttt{rpy2} setup on
     Linux/macOS/Windows with troubleshooting notes. \\
  15 & Tutorial Jupyter notebooks reproducing key textbook examples
     from Chapters~2, 5, and~6 of Maronna et al.\ (2019). \\
  16 & Final integration tests across all modules end-to-end.
     Contributor guide and README polish. \\
  \bottomrule
\end{tabularx}

\subsection{Phase 6: Stretch Wrappers, Native Python Track, and Final Submission (Weeks 17--22, Sep 21 -- Nov 2)}

\begin{tabularx}{\linewidth}{@{} l X @{}}
  \toprule
  \textbf{Week} & \textbf{Deliverables} \\
  \midrule
  17--18 & \textbf{Stretch wrappers (external packages):}
     \texttt{pense}, \texttt{GSE}, \texttt{TSGS} where time permits.
     Pytest suite. \\
  19--20 & \textbf{Native Python port and comparison (secondary track):}
     consolidate native reference code for selected functions where
     available; run side-by-side benchmarks vs.\ \texttt{rpy2} wrappers
     and R; document methodology, scope, and limits. \\
  21 & Final documentation review: ensure all wrappers are covered,
     all examples run, all tests pass. Address outstanding
     mentor feedback. \\
  22 & \textbf{Final submission:}
     Complete code review, final evaluation report, archive of all
     notebooks and benchmarks. \\
  \bottomrule
\end{tabularx}

%--------------------------------------------------------------------
% SECTION 10 -- RESULTING ROBSTATPY FEATURES
%--------------------------------------------------------------------
\section{Resulting RobstatPY Features}
\label{sec:resulting-robstatpy}

\noindent
\textbf{RobstatPY} is the outcome of the project: a pip-installable
\texttt{robstatpy} package that brings the in-scope \texttt{RobStatTM}
function set (Section~4) into Python with the same inputs, outputs, and
numerical behaviour as R through lightweight \texttt{rpy2} wrappers, so
users gain the full estimation capability of those methods without
reimplementing them.

\subsection{Validation, benchmarks, and documentation}

\begin{enumerate}[leftmargin=2em, itemsep=4pt]
  \item \textbf{Pytest test suite} achieving $\geq 90\%$ coverage
    across all wrapper functions, with field-by-field comparison
    against direct R calls confirming zero numerical difference for
    every output.

  \item \textbf{Comparative benchmark notebook} evaluating speed
    (native R vs.\ \texttt{rpy2} wrappers) and accuracy of the wrappers
    against R across all delivered functions, on both synthetic contaminated
    data and real book datasets (mineral, wine, bus).

  \item \textbf{Sphinx API documentation} published on ReadTheDocs,
    with docstrings modelled on the \texttt{RobStatTM} R man pages
    and runnable usage examples for each function.

  \item \textbf{Tutorial Jupyter notebooks} reproducing key textbook
    examples from Chapters~2, 5, and~6 of Maronna et al.\ (2019),
    demonstrating end-to-end workflows in Python.
\end{enumerate}

\subsection{External package wrappers (stretch)}

\noindent
\textbf{External package wrappers (first priority):}
\texttt{pense} (robust MM elastic net), \texttt{GSE} (generalised
S-estimator for missing data), and \texttt{TSGS} (two-step cell-wise
outlier estimator) --- recommended alongside \texttt{RobStatTM} functions
in the book but residing in separate R packages.

%--------------------------------------------------------------------
% SECTION 11 -- DISTRIBUTION
%--------------------------------------------------------------------
\section{Distribution and Packaging}

\texttt{robstatpy} will be distributed through two complementary
channels to maximise accessibility:

\subsection{GitHub Installation (Development and Bleeding-Edge)}

The package will be hosted on GitHub under an open-source license
(MIT or Apache~2.0, to be confirmed with mentors).
Users can install directly from the repository:

\begin{lstlisting}[style=python, title={\footnotesize GitHub installation}, numbers=none]
pip install git+https://github.com/Aakarsh751/robstatpy.git
\end{lstlisting}

\noindent
This channel provides access to the latest development version and
makes it straightforward for contributors to fork, modify, and submit
pull requests.
The repository will include:
\begin{itemize}[leftmargin=1.5em, itemsep=2pt]
  \item A \texttt{pyproject.toml} with full metadata and dependency
    specifications (\texttt{rpy2 >= 3.5}, \texttt{numpy},
    \texttt{pandas}, \texttt{scipy}, \texttt{matplotlib}).
  \item GitHub Actions CI running the full \texttt{pytest} suite on
    Linux, macOS, and Windows against multiple R versions (4.3, 4.4,
    4.5).
  \item A \texttt{CONTRIBUTING.md} with guidelines for adding new
    wrappers or submitting pure-Python ports.
  \item Tagged releases corresponding to each PyPI publication.
\end{itemize}

\subsection{PyPI Publication (Stable Releases)}

Stable, versioned releases will be published on the Python Package Index
(PyPI), enabling standard installation:

\begin{lstlisting}[style=python, title={\footnotesize PyPI installation}, numbers=none]
pip install robstatpy
\end{lstlisting}

\noindent
The packaging workflow uses:
\begin{itemize}[leftmargin=1.5em, itemsep=2pt]
  \item \texttt{build} to produce source distributions (\texttt{sdist})
    and pure-Python wheels.
  \item \texttt{twine} for secure upload to PyPI.
  \item Semantic versioning (e.g., \texttt{0.1.0} for the initial GSoC
    release, \texttt{0.2.0} after stretch wrappers are added).
  \item A GitHub Actions release workflow that automatically publishes
    to PyPI when a new git tag is pushed.
\end{itemize}

\noindent
\textbf{R dependency handling:}\/
Since \texttt{robstatpy} requires a working R installation with
\texttt{RobStatTM}, the package will:
\begin{enumerate}[label=(\arabic*), leftmargin=2.2em, itemsep=2pt, topsep=2pt]
  \item detect \texttt{R\_HOME} automatically at import time;
  \item provide a clear \texttt{ImportError} message with installation
    instructions if R is not found;
  \item include a \texttt{robstatpy.check\_setup()} utility that validates
    the R environment and prints version information.
\end{enumerate}
\noindent
The installation guide will cover R setup on all three major platforms:
Linux (\texttt{apt} or \texttt{conda}), macOS (Homebrew), and Windows
(the CRAN installer).

%--------------------------------------------------------------------
% SECTION 12 -- RISKS
%--------------------------------------------------------------------
\section{Risks and Mitigation Strategies}

\renewcommand{\arraystretch}{1.3}
\begin{tabularx}{\linewidth}{@{} >{\raggedright\arraybackslash}p{3.5cm}
                                  >{\raggedright\arraybackslash}X
                                  >{\raggedright\arraybackslash}X @{}}
  \toprule
  \textbf{Risk} & \textbf{Impact} & \textbf{Mitigation} \\
  \midrule
  \texttt{rpy2} build failures on Windows/macOS &
  Delays CI and user installation &
  Provide pre-built wheels where possible; document Conda/conda-forge
  installs and Docker-based workflows as fallbacks; pin supported R
  versions in CI \\[4pt]

  R or \texttt{RobStatTM} version incompatibility &
  Wrapper outputs may differ across R versions; breaking API changes
  could invalidate tests &
  Pin minimum R version (4.3+) and \texttt{RobStatTM} version;
  test CI matrix against R 4.3, 4.4, 4.5; add version check at
  import time \\[4pt]

  Complex rpy2 type conversion edge cases &
  Some \texttt{RobStatTM} functions return nested R lists or
  named vectors that do not map cleanly to Python objects &
  Each wrapper undergoes field-by-field validation against direct
  R output; conversion helpers (\texttt{rx}, \texttt{r2py}) are
  tested independently before use \\[4pt]

  Scope of 11 core functions proving too ambitious &
  Later phases (multivariate, PCA) could be under-delivered &
  Phases are ordered by dependency; univariate and regression
  wrappers (Phases 1--2) are self-contained and sufficient for
  a strong midterm; stretch functions absorb any schedule risk \\[4pt]

  Mentor availability during summer &
  Reduced feedback cadence &
  Three mentors provide redundancy; asynchronous communication
  via GitHub issues; weekly written status updates regardless \\
  \bottomrule
\end{tabularx}
\renewcommand{\arraystretch}{1}

%--------------------------------------------------------------------
% SECTION 13 -- REFERENCES
%--------------------------------------------------------------------
\section{References}

\begin{enumerate}[leftmargin=2em, itemsep=3pt, label={[\arabic*]}]
  \item R.~A.~Maronna, R.~D.~Martin, V.~J.~Yohai, and
    M.~Salibian-Barrera.
    \textit{Robust Statistics: Theory and Methods (with R)}, 2nd ed.
    Wiley, 2019.
  \item M.~Salibian-Barrera, G.~Willems, and R.~Zamar.
    ``The fast-$\tau$ estimator for regression.''
    \textit{Journal of Computational and Graphical Statistics}, 2008.
  \item V.~J.~Yohai.
    ``High breakdown-point and high efficiency robust estimates for
    regression.''
    \textit{The Annals of Statistics}, 15(2):642--656, 1987.
  \item D.~L.~Rocke.
    ``Robustness properties of S-estimators of multivariate location and
    shape in high dimension.''
    \textit{The Annals of Statistics}, 24(3):1327--1345, 1996.
  \item L.~Gautier et al.
    \texttt{rpy2}: R in Python.
    \url{https://rpy2.github.io/}, 2024.
  \item R.~D.~Martin.
    \texttt{PCRA}: Portfolio Construction and Risk Analysis.
    CRAN, 2023.
\end{enumerate}

\end{document}
